%! Solving by Square Roots & Completing the Square — Unit 3, Lesson 3.3 — Guided Notes (Answer Key)
\documentclass[10pt]{article}
\usepackage{algebra2-article}
\usepackage{algebra2-key}   % pulls in -boxes; do NOT also load -boxes
\newcommand{\TallMath}[1]{$\displaystyle #1\rule[-1.4em]{0pt}{3.2em}$}
\newcommand{\lgrid}[4]{%
  \draw[step=1, linegray!40, very thin] (#1,#3) grid (#2,#4);
  \draw[->, semithick, charcoal] (#1,0)--(#2,0) node[below right, font=\scriptsize]{$x$};
  \draw[->, semithick, charcoal] (0,#3)--(0,#4) node[left, font=\scriptsize]{$y$};}
% Vocab answer for the key: bold forest term, then the filled definition on a write-line.
\newcommand{\vocabans}[2]{%
  \noindent\textbf{\textcolor{forest}{#1:}}\\[1pt]\ansline{#2}}

\begin{document}

\pageheader{Unit 3, Lesson 3.3}{Guided Notes --- Answer Key}
\namedateperiod

\vspace{0.10in}

\begin{objectivebox}
\textbf{By the end of this lesson, I will be able to\dots}
\begin{itemize}\setlength{\itemsep}{2pt}
  \item solve $x^2=k$ and $(x-h)^2=k$ with the \emph{Square Root Property}, keeping the $\pm$ and simplifying radicals;
  \item \emph{complete the square} to solve any quadratic with $a=1$, even when it will not factor;
  \item rewrite $x^2+bx+c$ in \emph{vertex form} $(x-h)^2+k$ and choose the smartest solving method.
\end{itemize}
\end{objectivebox}

\vspace{0.06in}

\begin{vocabbox}
\small
Fill in each term as we build it together.
\vspace{2pt}
\vocabans{Square Root Property}{if $x^2=k$ with $k\ge0$, then $x=\pm\sqrt{k}$ --- take the square root of both sides, keeping $\pm$}
\vocabans{$\pm$ (plus or minus)}{records both square roots of a number; a squared equation has two solutions}
\vocabans{Simplest radical form}{a radical with no perfect-square factor left inside, e.g. $\sqrt{12}=2\sqrt{3}$}
\vocabans{Completing the square}{adding $(b/2)^2$ to $x^2+bx$ to build the perfect-square trinomial $(x+\tfrac{b}{2})^2$}
\vocabans{Vertex form}{$a(x-h)^2+k$, vertex $(h,k)$; completing the square converts standard form into it}
\end{vocabbox}

\begin{hookbox}
Two quadratics that \textbf{do not factor} over the integers:
\[ x^2-6=0 \qquad\text{and}\qquad x^2+6x+1=0. \]
\begin{itemize}\setlength{\itemsep}{1pt}
  \item Try to factor $x^2-6$: two integers with product $-6$ and sum $0$? \ansline{None exist --- $\sqrt6$ is irrational, so it won't factor over the integers.}
  \item Does that mean these have no solutions? (Their parabolas cross the $x$-axis!) \ansline{No --- the roots are real but irrational; we just need a non-factoring method to find them.}
\end{itemize}
Today we solve them exactly --- with two tools that don't need factoring.
\end{hookbox}

\begin{notesbox}{1. The Square Root Property --- undo the square}
``What squares to $9$?'' Both $3$ \emph{and} $-3$. So a squared equation has \textbf{two} solutions:
\[ \text{if } x^2=k \ (k\ge0), \text{ then } x=\textcolor{keyred}{\mathbf{\pm\sqrt{k}}}. \]
The $\pm$ is not optional --- forgetting it loses a root.
\begin{center}
\renewcommand{\arraystretch}{1.5}
\begin{tabularx}{\linewidth}{@{}X X@{}}
\textbf{Isolate the square first, then root.} & \textbf{Simplify the radical.} \\
$2x^2-50=0 \Rightarrow x^2=\textcolor{keyred}{\mathbf{25}} \Rightarrow x=\textcolor{keyred}{\mathbf{\pm5}}$ &
$3x^2=36 \Rightarrow x^2=\textcolor{keyred}{\mathbf{12}} \Rightarrow x=\textcolor{keyred}{\mathbf{\pm2\sqrt3}}$ \\
\end{tabularx}
\end{center}
Irrational is fine: $x^2-6=0 \Rightarrow x=\textcolor{keyred}{\mathbf{\pm\sqrt6}}$.

\vspace{3pt}
\fbox{\parbox{0.96\linewidth}{\small\textbf{The wall (preview):} $x^2=-4$ has \textbf{no real} solution ---
no real number squares to a negative. Lesson~3.4 builds the numbers that fix this. Today every radicand
stays $\ge0$.}}
\end{notesbox}

\begin{notesbox}{2. When the square is a binomial}
The rule is the same when a whole \emph{binomial} is squared --- root both sides, then undo the shift.
\[ (x-3)^2=16 \Rightarrow x-3=\textcolor{keyred}{\mathbf{\pm4}} \Rightarrow x=\textcolor{keyred}{\mathbf{7}} \ \text{or} \ x=\textcolor{keyred}{\mathbf{-1}}. \]
\[ (x+2)^2=12 \Rightarrow x+2=\textcolor{keyred}{\mathbf{\pm2\sqrt3}} \Rightarrow x=\textcolor{keyred}{\mathbf{-2\pm2\sqrt3}}. \]
This $(x-h)^2=k$ shape is exactly what the next method \emph{manufactures}.
\end{notesbox}

\begin{notesbox}{3. Completing the square ($a=1$)}
When a $bx$ term blocks the square-root move, \emph{build} a perfect square. From the warm-up,
$x^2+bx+\left(\tfrac{b}{2}\right)^2=\left(x+\tfrac{b}{2}\right)^2$, so add $\left(\tfrac{b}{2}\right)^2$ to
\textbf{both} sides.
\smallskip

\textbf{Solve} $x^2+6x+1=0$:
\begin{enumerate}[leftmargin=1.7em, itemsep=2pt]
  \item Move the constant: $x^2+6x=\textcolor{keyred}{\mathbf{-1}}$.
  \item Half of $6$ is $\textcolor{keyred}{\mathbf{3}}$; square it: $\textcolor{keyred}{\mathbf{9}}$. Add to both sides: $x^2+6x+9=\textcolor{keyred}{\mathbf{8}}$.
  \item Factor the perfect square: $(x+\textcolor{keyred}{\mathbf{3}})^2=8$.
  \item Root and solve: $x+3=\textcolor{keyred}{\mathbf{\pm2\sqrt2}} \Rightarrow x=\textcolor{keyred}{\mathbf{-3\pm2\sqrt2}}$.
\end{enumerate}
So $x=-3\pm2\sqrt2$ --- the two irrational $x$-intercepts of the parabola.
\end{notesbox}

\begin{notesbox}{4. Completing the square $=$ vertex form; choosing a method}
The \emph{same} move rewrites a quadratic in vertex form (Lesson~3.1) --- so it also finds the vertex.
\begin{center}
\begin{minipage}{0.56\linewidth}
\textbf{Rewrite the anchor} $x^2-2x-3$ in vertex form:
\[ x^2-2x-3=(x-\textcolor{keyred}{\mathbf{1}})^2-\textcolor{keyred}{\mathbf{4}}. \]
(Half of $-2$ is $-1$; add and subtract $1$.) Vertex $(\,\textcolor{keyred}{\mathbf{1}},\,\textcolor{keyred}{\mathbf{-4}}\,)$ --- matching 3.0--3.1.

\vspace{4pt}
\renewcommand{\arraystretch}{1.3}
\begin{tabular}{|l|l|}
\hline
\textbf{If the quadratic\dots} & \textbf{use\dots} \\ \hline
factors nicely & factoring (3.2) \\ \hline
has no $bx$ term & square roots \\ \hline
won't factor & completing the square \\ \hline
\end{tabular}
\end{minipage}%
\begin{minipage}{0.42\linewidth}\centering
\begin{tikzpicture}[scale=0.42]
  \lgrid{-3}{5}{-5}{3}
  \begin{scope}\clip (-3,-5) rectangle (5,3);
    \draw[very thick, forest, domain=-1.4:3.4, samples=75, smooth] plot (\x,{(\x-1)*(\x-1)-4});
  \end{scope}
  \fill[charcoal] (-1,0) circle (3pt) node[above left, font=\scriptsize]{$-1$};
  \fill[charcoal] (3,0) circle (3pt) node[above right, font=\scriptsize]{$3$};
  \fill[charcoal] (1,-4) circle (3pt) node[below, font=\scriptsize]{$(1,-4)$};
\end{tikzpicture}
\end{minipage}
\end{center}
\end{notesbox}

\begin{practicebox}
\textbf{Solve. Show the isolated square (or completed square), then both roots in simplest form.}
\begin{enumerate}[leftmargin=1.9em, itemsep=6pt]
  \item $x^2-20=0 \Rightarrow x=\textcolor{keyred}{\mathbf{\pm2\sqrt5}}$
  \item $(x+1)^2=9 \Rightarrow x=\textcolor{keyred}{\mathbf{2,\,-4}}$
  \item $x^2+4x-6=0$ (complete the square) $\Rightarrow x=\textcolor{keyred}{\mathbf{-2\pm\sqrt{10}}}$
\end{enumerate}
\end{practicebox}

\begin{teachernote}
The two engines: the Square Root Property (\S1--2, always with $\pm$ and simplest radical form) and
completing the square (\S3, the $(b/2)^2$ recipe from the warm-up). \S4 is the payoff --- the \emph{same}
move gives vertex form, so it ties back to the anchor $x^2-2x-3=(x-1)^2-4$, vertex $(1,-4)$, roots $-1,3$.
Pre-empt the four recurring slips: dropping the $\pm$; rooting before isolating (\S1 ``isolate first'');
using $b$ instead of $b/2$; and unsimplified radicals ($\sqrt{20}$ vs.\ $2\sqrt5$ in practice~1). The
$x^2=-4$ ``wall'' is the deliberate bridge to Lesson~3.4 --- name it as ``no \emph{real} solution,'' not
``no solution.''
\end{teachernote}

\end{document}
